\documentclass[11pt,a4paper]{article}
\usepackage[utf8]{inputenc}
\usepackage[T1]{fontenc}
\renewcommand{\abstractname}{Összefoglaló}
\renewcommand{\contentsname}{Tartalomjegyzék}
\usepackage{lmodern}
\usepackage[margin=2.2cm]{geometry}
\usepackage{amsmath,amssymb,amsthm}
\usepackage{booktabs,array,longtable}
\usepackage{xcolor}
\usepackage{enumitem}
\usepackage[colorlinks=true,linkcolor=blue!60!black,citecolor=blue!60!black,urlcolor=blue!60!black]{hyperref}
\setlength{\emergencystretch}{3em}
\setlength{\parskip}{0.45em}
\setlength{\parindent}{0pt}

\title{\textbf{A Tau Core kutatási program fejl\H{o}déstörténete}\\[0.4em]
\large Az alapötlett\H{o}l (\texttt{GalaxyTest} / \texttt{SRTM} / \texttt{MTW}) a \texttt{TPG}-n és a \texttt{HDDA--DTL} falszifikációs keretrendszeren át az egyesített $K_4$ szimplex elméletig (\texttt{v27})}
\author{Olcsák József\\{\small Független kutató, Budapest (\texttt{jolcsak@gmail.com})}\\{\small ORCID: \href{https://orcid.org/0009-0003-7311-6522}{0009-0003-7311-6522}}}
\date{2026. november 15. (\texttt{v27} kanonikus kiadás -- 49 PDF, 33/33 ellen\H{o}rz\H{o} kapu)}

\begin{document}
\maketitle

\begin{abstract}
\noindent
Ez a dokumentum a \textbf{Tau Core} kutatási program teljes történeti, fogalmi és matematikai ívét foglalja össze a legels\H{o} alapötlett\H{o}l a 2026. november 15-i \texttt{v25}-ös egyesített kiadásig. A program abból a fizikai alapfeltevésb\H{o}l indult ki, hogy az id\H{o} ($\tau$) nem passzív háttérkoordináta, hanem bels\H{o} spektrummal és véges információs kapacitással rendelkez\H{o} fizikai szubsztrátum, amelynek megfigyel\H{o}i levetítéséb\H{o}l (\emph{$\tau$-projekció}) a galaktikus forgási görbék, a gravitáció és a kvantumos jelenségek származnak. Nyomon követjük, hogyan kényszerítették ki a szigorú önfalszifikációs kapuk a modell egymást követ\H{o} fejl\H{o}dési lépcs\H{o}it:
(1) a 175 \textsc{sparc} galaxison tesztelt korai \texttt{GalaxyTest} $\to$ \texttt{SRTM} $\to$ \texttt{MTW} (\texttt{V1--V6}) fázist;
(2) a \textbf{\texttt{TPG} ($\tau$-Projection Gravity)} és az 5-dimenziós \textsc{rmond} kompatibilitási hidat;
(3) a Hollinger-féle szuperhalmaz-kapu (\texttt{18/18 FAIL}) által kikényszerített \textbf{\texttt{HDDA--DTL}} (\emph{History-Dependent Deformation Algebra / Dynamic Temporal Layers}) korszakot és a 2026. május 8-i ontológiai fordulatot;
(4) a \textbf{\texttt{Tau Core}} és a \textbf{\texttt{TPST}} kiválását, valamint az Empirikus Sorozatot (\texttt{Papers 1--19});
(5) a Római Alapozó Sorozatot (\texttt{Papers I--VII-C}) és az \texttt{F16--F20} matematikai blokkoló-feloldást; végül
(6) a jelenlegi, \textbf{$k=0$ szabad paraméteres, hárompilléres $K_4$ tetraéderes szimplex szintézist} (\texttt{C1--C4}, \texttt{Papers 20--22}), amely egyetlen közös geometriai forrásból írja le a mikrovilágot (kvantummechanika, elektromágnesesség, \textrm{Spin}(10) részecskefizika, UV-véges kvantumgravitáció) és a makrovilágot (fekete lyukak, galaxisok, kozmikus id\H{o}térkép és tágulás).
\end{abstract}

\tableofcontents
\bigskip
\hrule
\bigskip

\section{Az alapötlet (0. fázis): Az id\H{o} mint els\H{o}dleges fizikai szubsztrátum}

A modern elméleti fizika két legnagyobb feszültsége egyrészt a kozmológiai sötét szektor (a galaxisok forgási görbéihez és a lencséz\H{o} tömeghez posztulált $\sim 85\%$ láthatatlan sötét anyag), másrészt az id\H{o} kett\H{o}s szerepe (az általános relativitáselméletben dinamikus sajátid\H{o}, a kvantummechanikában küls\H{o} paraméter).

A kutatási program alapötlete egyetlen fizikailag szemléletes kérdésb\H{o}l született:
\begin{quote}
\emph{Mi van akkor, ha a galaxisok küls\H{o} tartományaiban tapasztalt többlet-gravitációt nem láthatatlan sötét anyag részecskék okozzák, hanem az, hogy az id\H{o} ($\tau$) egy mélyebb, többréteg\H{u} spektrum, amelynek a küls\H{o} megfigyel\H{o} csak a levetített (integrált) hatását érzékeli?}
\end{quote}
Ha a lokális newtoni gravitációs gyorsulás ($a_N$) egy kritikus kozmikus küszöb ($a_0 \approx 1{,}2\times 10^{-10}\;\mathrm{m/s^2}$) alá csökken, a $\tau$-projekció logaritmikus többlet-er\H{o}t hoz létre új részecskék bevezetése nélkül.

Ebb\H{o}l a kiinduló kérdésb\H{o}l három, egymásra épül\H{o} fizikai alapintuíció bontakozott ki, amelyek kés\H{o}bb a teljes matematikai elmélet tartóoszlopaivá váltak:
\begin{enumerate}[leftmargin=1.5em]
\item \textbf{A teljes id\H{o}vonalra kiterjed\H{o} galaktikus morfológia (múlt, jelen és jöv\H{o} együttes hatása):} Az a felismerés, hogy egy galaxis forgását nem lehet egyetlen pillanatnyi 3D-s felvételbe zárva megérteni, hanem a jelen mellett a múlt és a jöv\H{o} is egyszerre hat a forgására --- vagyis a galaxis morfológiáját a teljes id\H{o}vonalra kiterítve kell ábrázolni. A kés\H{o}bbi matematikai levezetésben pontosan ez öltött testet egyrészt a Cartan--Kodama küls\H{o} derivált Faraday-id\H{o}indukciós tagjában ($\partial_t K_r = -\dot{H}r$), amely az $a_0^{\mathrm{KH}}(z) = \frac{5\sqrt{3}}{96}cH(z)|1 - q(z)|$ képletben összekapcsolja a múltbeli fékez\H{o}dést ($q>0$) és a jöv\H{o}beli gyorsulást ($q<0$), létrehozva a $z_{\mathrm{dip}} \equiv z_{\mathrm{acc}} = 0{,}6316$-os ($q=0$) galaktikus völgyet; másrészt a kétvég\H{u} megfigyel\H{o}--forrás oksági folyosóban ($T_{P|OS}$) és a $K_4$ $(2+2)$-es kauzális vágásán fennálló Osterwalder--Schrader reflexiós pozitivitásban.
\item \textbf{Az atemporális Morfológiai Testet ($M_\tau^\star$) létrehozó id\H{o} el\H{o}tti kölcsönhatás (\texttt{MRC-DEF1/DEF2}):} Annak végiggondolása, hogy ha ez a Morfológiai Test már magában foglalja a teljes id\H{o}vonalat, akkor az a kölcsönhatás, amely kialakítja, nem zajlódhat le a koordinátaid\H{o}ben. Így született meg a tér és id\H{o} el\H{o}tti, $S_4$-szimmetrikus \textbf{Alapmez\H{o} ($B_\tau^{\mathrm{base}}$)} és az \emph{egyetlen univerzum-szint\H{u}}, szimmetriasért\H{o}, idempotens \textbf{Csíra-projektor ($s_U$, $r_{\mathrm{seed}} = 1/4$)} stacionárius variációs egyensúlya ($K_B m = j_s$), amelynek atemporális válasza maga az \emph{egyetlen univerzum-szint\H{u}} \textbf{Morfológiai Test} ($[M_\tau^\star] = [\mathcal{F}_\tau(B_\tau^{\mathrm{base}}, s_U, \mathcal{K})]_\sim$). Az \texttt{MRC-DEF2} szabály szerint sem $B_\tau^{\mathrm{base}}$, sem $s_U$ nem 4D-s barionos anyag vagy térid\H{o}beli tárgy; a galaxisok, kozmikus üregek, atomok, fotonok és megfigyel\H{o}k az egyetlen $M_\tau^\star$ testb\H{o}l leszálló alrendszer-leírók ($I_{A|S}(o)$), soha nem külön-külön csíráztatott szül\H{o}-testek.
\item \textbf{Anyagtól és id\H{o}t\H{o}l független közös univerzális nézet:} Az a felismerés, hogy az atemporális Morfológiai Test ($M_\tau^\star$) egy anyagtól, skálától és id\H{o}t\H{o}l független közös geometriai nézet: amit a bels\H{o} megfigyel\H{o} ,,anyagnak'' ($O_M$), ,,id\H{o}nek'' ($O_T[M_\tau^\star](\mathbf{x},z) = \mathrm{d}\tau/\mathrm{d}t(\mathbf{x},z)$ 4D-s terminális óraráta-kiolvasásnak), ,,kvantumos hullámfüggvénynek'' ($O_Q$), ,,gravitációnak'' ($O_G$) vagy ,,látszólagos sötét energiának'' érzékel, az mind ugyanannak az egyetlen közös $K_4$ Morfológiai Testnek ($N=4 \implies u = \frac{\sqrt{3}}{24}$, $\alpha_{\mathrm{OCSR}} = 5u = \frac{5\sqrt{3}}{24}$) a különböz\H{o} lépték\H{u} 4D-s kiolvasása a rendezett megkülönböztethet\H{o}ségi hordozón ($(A_{OS}, D_{OS}, \le_{OS})$) keresztül.
\end{enumerate}

\section{I. Korszak: \texttt{GalaxyTest} $\to$ \texttt{SRTM} $\to$ \texttt{MTW} (2026. április 18--28.)}

Az alapötlet els\H{o} tesztelése egy determinisztikus C\# / .NET számítási környezetben (\texttt{GalaxyTest}) indult el a \textsc{sparc} adatbázis galaxisain:

\begin{enumerate}[leftmargin=1.5em]
\item \textbf{\texttt{Rev 1} (2026.04.18) -- A 7 galaxisos prototípus és a \texttt{V1--V5} korai modellváltozatok korlátai:}
Az els\H{o} öt modellváltozat (\texttt{V1--V5}), amelyek a korong skálahosszától ($R_{\mathrm{disk}}$) és s\H{u}r\H{u}ség-modulációtól függtek, 7 galaxison ígéretesnek t\H{u}ntek ($\alpha = 0{,}405$), de a teljes 175 \textsc{sparc} galaxisra kiterjesztve nagyobb szórást adtak (\textsc{rmse} $25\text{--}67\;\mathrm{km/s}$). Ezeket a korai közelítéseket a projekt dokumentálta és továbbfejlesztette.
\item \textbf{\texttt{Rev 2} (2026.04.19) -- A \texttt{V6} logaritmikus gyorsulás-válasz eredménye 175 galaxison:}
A skálahossz-függés helyett bevezetett \textbf{\texttt{V6} képlet}:
\begin{equation}
F(a_N) = 1 + \alpha\,\ln\!\left(1 + \frac{a_0}{a_N}\right), \qquad \alpha = 0{,}360
\end{equation}
a teljes 175 \textsc{sparc} galaxison $16{,}62\;\mathrm{km/s}$ átlagos hibát (\textsc{rmse}) ért el (a \textsc{mond} $17{,}10\;\mathrm{km/s}$ értékével szemben, $0{,}97\times$ \textsc{mond} arány), különösen jól teljesítve a Tejútrendszeren ($0{,}90\times$ \textsc{mond}) és a kis felületi fényesség\H{u} (\textsc{lsb}) spirálgalaxisokon.
\item \textbf{Névváltás \texttt{SRTM}-r\H{o}l \texttt{MTW}-re (2026.04.25):}
A legels\H{o} \textbf{\texttt{SRTM}} (\emph{Scale-Relativistic Time Metric}) munkanevet elhagyták, mert a Nottale-féle skála-relativitás félrevezet\H{o} volt a tényleges matematikai szerkezethez (logaritmikus súlyozású $\tau$-integrál) képest. Az új név \textbf{\texttt{MTW} (\emph{Multi-Time Window Gravity})} lett (\texttt{jolcsak/GalaxyTest} $\to$ \texttt{jolcsak/MTW}).
\item \textbf{AQUAL Lagrange-függvény és az els\H{o} kozmológiai kapcsolat (\texttt{Rev 3--6}):}
Elkészült a nem-relativisztikus \textsc{aqual} Lagrange-levezetés (\texttt{TpgAqualLagrangian.cs}, $3\times 10^{-16}$ gépi pontossággal), a 123 regressziós teszthorgony, valamint az a felismerés, hogy az empirikusan illesztett $\alpha \approx 0{,}360$ szorosan kapcsolódik a Hubble-állandóhoz:
\begin{equation}
\alpha \approx \frac{2 a_0}{c H_0} \approx 0{,}366.
\end{equation}
\end{enumerate}

\section{II. Korszak: \texttt{TPG} ($\tau$-Projection Gravity) és az 5D \textsc{rmond}-híd (2026. április 29. -- május 4.)}

\subsection{Miért lett az \texttt{MTW}-b\H{o}l \texttt{TPG}?}
2026. április 29-én a projektet \textbf{\texttt{TPG} (\emph{$\tau$-Projection Gravity}, $\tau$-Projekciós Gravitáció)} névre keresztelték át (\texttt{jolcsak/MTW} $\to$ \texttt{jolcsak/TPG}), hogy elkerüljék a névütközést Misner, Thorne és Wheeler klasszikus \emph{Gravitation} (1973, ``\textsc{mtw}'') tankönyvével, és pontosan megnevezzék az 5. axiómát: \emph{a megfigyel\H{o} nem csatolódik közvetlenül a rejtett $\tau$ dimenzióhoz, csak a $\tau$-integrált mennyiségek mérhet\H{o}k}.

\subsection{Az 5D \texttt{TPG--\textsc{rmond}} kompatibilitási híd}
A \textsc{gw170817} gravitációshullám-mérés ($|c_g - c|/c < 10^{-15}$) kizárta a korábbi relativisztikus \textsc{mond}-kísérleteket (\textsc{teves}), de Skordis és Z\l{}o\'snik (2020, \texttt{SZ20}) megalkotta a kovariáns \textsc{rmond} elméletet egy $\varphi$ skalármez\H{o}vel és egy $A^\mu$ id\H{o}szer\H{u} egységvektorral, $\sim 10$ szabad paraméterrel. A \texttt{TPG} híd bebizonyította, hogy egy 5-dimenziós ($\mathbb{R}^4 \times \mathbb{R}_\tau$) szeparábilis mez\H{o}feltevés és a
\begin{equation}
\mathcal{O}_{\mathrm{obs}}(x) = \int_0^\infty K_q(\tau)\,\mathcal{O}(x,\tau)\,d\tau, \qquad K_q(\tau) = \frac{\alpha}{(1+\tau)^{1+q}}, \qquad q_{\mathrm{ch}} = \alpha\,k_{\mathrm{ch}}^{1/(k_{\mathrm{ch}}-1)}
\end{equation}
lezárási mag ($k_v = 2 \implies q_v = 2\alpha$, $k_s = 3 \implies q_s = \alpha\sqrt{3}$) $\tau$-integrálásával az \texttt{SZ20 \textsc{rmond}} hatásfüggvény \textbf{mechanikusan el\H{o}áll egyetlen $\alpha$ paraméterb\H{o}l}, $\alpha$-független $c_{\mathrm{mix}} = 1/(4+\sqrt{3})$ keveredési együtthatóval.

\subsection{A falszifikációs fal, amely kikényszerítette a továbblépést}
A \texttt{TPG} szigorú önellen\H{o}rz\H{o} kapui három kritikus hiányosságot tártak fel:
\begin{enumerate}[leftmargin=1.5em]
\item \textbf{Hollinger lokális szuperhalmaz skalár-kapu (\texttt{18/18 FAIL}):} Az árnyékolatlan, pusztán a lokális $a_N$ gyorsulástól függ\H{o} logaritmikus \texttt{TPG/v6} válasz 18-ból 18 szuperhalmaz-soron megbukott, míg egy környezeti s\H{u}r\H{u}séggel kapuzott (\emph{density-gated}) változat átment.
\item \textbf{Galaxishalmaz-lencsézés (Coma-halmaz, \texttt{J7}):} A naiv 1-paraméteres \texttt{TPG} a Coma-halmaz lencséz\H{o} tömegének csak a $65{,}6\%$-át adta vissza.
\item \textbf{5D szeparábilis törékenység:} A 21 \texttt{SymPy}-kapus audit kimutatta, hogy az 5D-s szeparábilis feltevések (\texttt{A1--A3}) sérülékenyek, ha a $\tau$-szeletek között derivált-csatolás lép fel.
\end{enumerate}

\section{III. Korszak: \texttt{HDDA--DTL} és a 2026. május 8-i ontológiai fordulat}

A Hollinger-kapu \texttt{18/18}-as bukása bizonyította be, hogy a gravitációs válasz nem lehet pusztán a pillanatnyi lokális gyorsulás függvénye: figyelembe kell venni a rendszer \textbf{környezeti, morfológiai és terhelési el\H{o}történetét}. Így jött létre 2026. május 5-én a \textbf{\texttt{HDDA--DTL}} keretrendszer (\texttt{jolcsak/HDDA-DTL}):

\begin{itemize}[leftmargin=1.5em]
\item \textbf{\texttt{HDDA} (\emph{History-Dependent Deformation Algebra} / El\H{o}történet-függ\H{o} Deformációs Algebra):}
Egy $\mathcal{A} = (S, \Omega, H)$ algebrai hármas, ahol $S$ a hordozó állapottér, $\Omega$ a m\H{u}velet- és projektorcsalád, $H$ pedig a terhelési el\H{o}történet. Egy $x$ környezeti/forrás-esemény hatására maga a m\H{u}veletcsalád deformálódik:
\begin{equation}
D_x(\mathcal{A}) = \big(S,\; D_x(\Omega),\; H\cdot x\big).
\end{equation}
\item \textbf{\texttt{DTL} (\emph{Dynamic Temporal Layers} / Dinamikus Id\H{o}rétegek):}
A fizikai alkalmazási réteg, amely a logaritmikus választ véges kapacitású id\H{o}rétegek (\emph{finite-capacity loading}) telít\H{o}déseként modellezte.
\item \textbf{A 2026. május 8-i ontológiai fordulat (\emph{HDDA Tau-Space Compatibility Refactor Gate v1}):}
A bels\H{o} elméleti audit kimondta, hogy a szó szerinti ``egymásra pakolt diszkrét id\H{o}lapok'' (\emph{literal discrete time sheets / layer hydrodynamics}) nem vezethet\H{o}k le kovariáns hatásfüggvényb\H{o}l. A projekt ezért elvetette a naiv réteg-ontológiát, és minden fogalmat átfordított a \textbf{folytonos/operátoros $\tau$-állapottér ($\Psi(x,\tau)$) projekciós nyelvére}: az ``id\H{o}réteg'' projektált véges rangú altérré, a ``rétegkapacitás'' projektor-ranggá ($\mathrm{rank}(P_\tau)$), a ``rétegáram'' pedig a forráshatás Noether-áramává ($\delta S_{\mathrm{src}} \to J_\tau$) vált.
\end{itemize}

\section{IV. Korszak: A \texttt{Tau Core} kiválása, a \texttt{TPST} és az Empirikus Sorozat (\texttt{Papers 1--19})}

2026. május 12-én a \textbf{\texttt{Tau Core}} önálló programmá vált. A központi objektum a $\Psi(x,\tau)$ állapot lett, amelyb\H{o}l az öt megfigyelhet\H{o} fizikai szektor öt kanonikus projekcióként (\emph{readout}) áll el\H{o}:
\begin{equation}
P_{\mathrm{time}}[\Psi],\qquad P_{\mathrm{quantum}}[\Psi],\qquad P_{\mathrm{matter}}[\Psi],\qquad P_{\mathrm{radiation}}[\Psi],\qquad P_{\mathrm{gravity}}[\Psi].
\end{equation}

\begin{enumerate}[leftmargin=1.5em]
\item \textbf{\texttt{TPST} (\emph{Tau Projection-State Theory}):}
A fizika alatti tiszta matematikai szubsztrátum ($P_i : \mathcal{T} \to \mathcal{S}_i$), amely bebizonyította a szál-finomítási (\emph{fibre-refinement}), egyenl\H{o} rangú no-go, súlyozott Schur-komplementer és lokális Legendre/Bregman-redukciós ($\Gamma = -H^{-1}B = J_{\mathrm{OS}}$) tételeket.
\item \textbf{Az Empirikus és Protokoll Sorozat (\texttt{Papers 1--19}):}
\begin{itemize}
\item \emph{Galaktikus koherencia és morfológia (\texttt{Papers 1--3, 8, 9, 11}):} Vak A/B/C szerkezeti zavartsági audit a \textsc{sparc} galaxisokon (\texttt{Paper 1}, Zenodo \textsc{doi}), majd küls\H{o} optikai/\textsc{hi} morfológiából számított forgásigörbe-el\H{o}rejelzés (\texttt{Paper 8--9}).
\item \emph{Mikrovilág-el\H{o}futárok (\texttt{Paper 4} Higgs-modul, \texttt{Paper 15} Flavor-geometria):} A bels\H{o} $\tau$-sajátérték-probléma ($L_\tau\chi_n = \lambda_n\chi_n$) tömegskálákat, a $\tau$-bázis forgása mértékmez\H{o}ket, a $\nu_D = 3/10$ (\texttt{Branch A}) Higgs-lokalizáció pedig a Higgs-önkölcsönhatást ($I_4(3/10)=0{,}13376$) és a top Yukawa-csatolást ($0{,}96155$) generálja.
\item \emph{Kozmológia, lencsézés és 4D kovariancia (\texttt{Papers 5--7, 10, 13, 14, 16--19}):} Véges memóriájú kozmológia (\texttt{Paper 5}), Hubble Frontier Fields er\H{o}s lencsézés (\texttt{Papers 6--7}), Friedmann-audit (\texttt{Paper 10}), \textsc{cmb} sötétszektor-szükségszer\H{u}ség (\texttt{Paper 13}), Cosmic Web tanú-protokoll (\texttt{Paper 14}), valamint a teljes 4D kereszttartományú kovariancia-értékelés (\texttt{Papers 16--17}) és foton-lencsézési tesztek (\texttt{Papers 18--19}).
\end{itemize}
\end{enumerate}

\section{V. Korszak: A Római Alapozó Sorozat (\texttt{Papers I--VII-C}) és a \texttt{TD-0..TD-III} gerinc}

2026 augusztusa és októbere között a teljes elméleti építmény szigorú axióma--tétel--bizonyítás láncba rendez\H{o}dött a 11 Római Alapozó Tanulmányban és a 7 Technikai Levezetési Kötetben:
\begin{itemize}[leftmargin=1.5em]
\item \textbf{A \texttt{Carrier $\to$ Seed $\to$ Morphological Body} (\texttt{MRC-DEF1/DEF2}) architektúra (\texttt{Paper I}):} Elkülönítette a pregeometriai hordozót ($B_\tau^{\mathrm{base}}$), az olvasható terhelési magot ($r_{\mathrm{seed}}$) és a stabilizált morfológiai választestet ($M_\tau^\star$), valamint levezette a $\Gamma_{\mathrm{CLP3}}$ izometrikus kvantum--gravitációs transzportot ($A_* = \hbar$).
\item \textbf{Kvantum-, id\H{o}- és mértékelméleti lezárás (\texttt{Papers II--VII-C}):} \textsc{gns}-egyértelm\H{u}ség (\texttt{II}), fonat-holonómia (\texttt{III}), Hodge--Clifford közös fénykúp (\texttt{IV}), temporális leszállás $d\tau/dt$-re (\texttt{V}), Born-szabály és Uhlmann-h\H{u}ség (\texttt{VI}), valamint az egyesített Einstein--Maxwell--anyag végállomás (\texttt{VII-A}), a 256-dimenziós $\mathrm{CAR}_7$ algebra (\texttt{VII-B}) és a 15-skalár (\texttt{CHDF}) renormálási csoport (\texttt{VII-C}).
\item \textbf{Az \texttt{F16--F20} blokkoló-feloldási lánc (2026.09.29 -- 11.12):} Bebizonyította a kétlépcs\H{o}s Schur-dilatációs pozitivitást (\texttt{F18-T1}, szellemmentesség), az \texttt{OCSR-H1} koercivitást ($\lambda_{\min} = +0{,}15975 > 0$) és a differenciális óraráták egzakt skálainvarianciáját (\texttt{F20-T1}).
\end{itemize}

\section{VI. Korszak: A Hárompilléres $K_4$ Szintézis, a \texttt{C1--C4} Lezárások és a \texttt{Papers 20--22} (\texttt{v25})}

2026 novemberében a \texttt{Master Synthesis} bebizonyította, hogy minden korábbi 5D-s \texttt{TPG} segédfeltevés és \texttt{HDDA--DTL} segédváltozó feleslegessé válik, mert a teljes 4D-s egyesített elmélet \textbf{három egzakt matematikai pillérre redukálódik, $k = 0$ szabad paraméterrel}:

\begin{enumerate}[leftmargin=1.5em]
\item \textbf{I. Pillér -- A $K_4$ tetraéderes szimplex, $\mathrm{Cl}^+(1,9)\subset\mathrm{Spin}(10)$ és $r_{\mathrm{seed}} = 1/4$:}
A 4 csúcsú tetraéder ($K_4$) normált Laplace-spektruma és az egyetlen univerzális mag ($r_{\mathrm{seed}} = 1/4$) rögzíti az összes alapvet\H{o} csatolást:
\begin{equation}
\alpha_{\mathrm{OCSR}} = \frac{5\sqrt{3}}{24} \approx 0{,}360844, \qquad \alpha_{\mathrm{KH}} = \frac{5\sqrt{3}}{48} \approx 0{,}180422, \qquad \frac{\Omega_A}{\Omega_b} = \frac{213\sqrt{3}-20}{65} = 5{,}368105.
\end{equation}
\emph{Történeti csúcspont:} A $K_4$ szimplexb\H{o}l levezetett $\alpha_{\mathrm{OCSR}} = 5\sqrt{3}/24 = 0{,}360844$ els\H{o} elvekb\H{o}l megmagyarázza, miért kapott a legels\H{o} 2026. áprilisi \texttt{TPG/v6} empirikus \textsc{sparc}-illesztés $\alpha = 0{,}360$-at!
\item \textbf{II. Pillér -- Horizont-termodinamika és 60 dekádos Frullani/Cauchy-renormálás:}
Egyértelm\H{u}en rögzíti az egyetlen kanonikus karakterisztikus gyorsulásskálát ($a_0^{\mathrm{KH}}(z) = \frac{5\sqrt{3}}{96}cH(z)|1 - q(z)|$, $a_0^{\mathrm{KH}}(0) = 0{,}9015\times 10^{-10}\;\mathrm{m/s^2}$) és az effektív kozmológiai integrációs állandót ($\Lambda_{\mathrm{eff}} = 3\Omega_\Lambda H_0^2/c^2$).
\item \textbf{III. Pillér -- Egységes 4D kovariáns dilogaritmusos gravitációs hatásfüggvény ($W_{\mathrm{COCA}}(s)$):}
Zárt alakban állítja el\H{o} az általános relativitáselmélet er\H{o}s téri határértékét (Cassini-korlát), a galaktikus \textsc{rar}-t ($0{,}082\;\mathrm{dex}$), a \textsc{cmb} akusztikus csúcsait ($H_0^{\mathrm{Tau}} = 67{,}3388\;\mathrm{km/s/Mpc}$, \texttt{C2}), a galaxishalmazok tiszta $\ell_\tau(M_b,z)=\sqrt{GM_b/a_0^{\mathrm{KH}}(z)}$ magprofilját ($\chi^2_{\mathrm{XCOP}}=27{,}36$ zérus $\Lambda$CDM $c_{200}$-importtal), a galaxiscsoportok gázmegtartását (\texttt{C3}, medián arány $0{,}9762$) és a fénysebesség\H{u} gravitációs hullámokat ($Z_T = 1$).
\end{enumerate}

Erre a három pillérre épült fel a három legújabb szintézistanulmány (\texttt{v23 $\to$ v25}):
\begin{itemize}[leftmargin=1.5em]
\item \textbf{Paper 20 (\emph{The Cosmic Time-Landscape}):} 18\,983 valós \textsc{esa Euclid dr1} rekordból rekonstruált 3D kozmikus id\H{o}térkép ($O_T[M_\tau^\star](\mathbf{x},z) = \mathrm{d}\tau/\mathrm{d}t$) és az \texttt{F50} egyesített Galileo/PTB/Cassini/Euclid atomóra-audit.
\item \textbf{Paper 21 (\emph{Black Holes, the Big Bang, and Present Cosmic Expansion}):} A fekete lyukak ($\sigma_{\mathrm{flow}}=-1$) és a táguló Univerzum ($\sigma_{\mathrm{flow}}=+1$) közös Painlevé--Gullstrand id\H{o}geometriája, 55 nagyságrendes egyesített UV--IR $\beta_{\mathrm{master}}^2(r,z) = \frac{r_s r^2}{r^3+r_s\ell_*^2}F_{\mathrm{OD}}^2(\ell_\tau^2/r^2)$ mesterprofil, szingularitásmentes $K_4$ mag ($\ell_* = \sqrt{24\pi}\ell_P$), $11{,}28\,M_P$ hideg maradvány, és a $z_{\mathrm{acc}} = z_{\mathrm{dip}} = 0{,}6316$ kozmikus iker-egybeesés.
\item \textbf{Paper 22 (\emph{The Micro-World of Tau Core}) és az \texttt{F52--F55} kereszttartományú konzisztencia-, axiomatikus tömörítési és áramlástani lezárás (\texttt{Gates E27--E31}, \texttt{35/35 PASS}):} A kvantummechanika, elektromágnesesség ($\sin^2\theta_W(M_Z) = 0{,}231220$ az $M_{\mathrm{PS}} = 3{,}66\times 10^{11}\;\mathrm{GeV}$ Pati--Salam/skalár-15 küszöbsugáron), \textrm{Spin}(10) neutrínótömegek ($(0,\;8{,}83,\;50{,}36)\;\mathrm{meV}$), a $G_\nu = -3S + 8A$ generátorból származó királis $CP$-szimmetriasértés ($\delta_{\mathrm{CP}} \approx -90^\circ$, $J_{\mathrm{PMNS}} \approx -0{,}0337$, miközben $\bar{\theta}_{\mathrm{QCD}} = 0$ axionok nélkül) és a 15 csatornás UV-véges kvantumgravitáció egyesítése a kozmológiával a közös $u = \frac{\sqrt{3}}{24}$ generátoron és az L1/L2 \H{u}rbeli nyeregponti $3{,}479\times$ \textsc{bmv} kvantum-összefonódási el\H{o}rejelzésen keresztül (\texttt{F51 / Gate E27}); kiegészítve a 4-utas, Planck-mentes $H_0 = 67{,}337 \pm 0{,}040\;\mathrm{km/s/Mpc}$ ($100\theta_*^{\mathrm{vak}} = 1{,}04104$, $-0{,}19\sigma$) konkordanciával, a 18 diszjunkt kis tömeg\H{u} galaxiscsoport küls\H{o} tesztjével ($\mathrm{medián} = 0{,}9762$, $\chi^2/N = 0{,}106$, \texttt{F52 / Gate E28}), a tartományonkénti homogén Likelihood-összehasonlítással ($\Delta_{\mathrm{OOS}}(2\mathrm{p}-K_4)=+161{,}25$, 1p-$u$: $+56{,}19$, \texttt{F53 / Gate E29}), az egyetlen topológiai egész számra ($N=4 \implies u=\frac{\sqrt{N-1}}{N!}=\frac{\sqrt{3}}{24}$, Tier-1/Tier-2 eredetszétválasztással) és egysoros $F_{\mathrm{OD}}^2(s)$ variációs magra épül\H{o} axiomatikus tömörítéssel (\texttt{F54 / Gate E30}), valamint a fekete lyukak bels\H{o} terének ($\sigma_{\mathrm{shear}}\equiv 0$ nyírásmentes de Sitter-mag, $M_{\mathrm{rem}}=11{,}28\,M_P$) és a Nagy Visszapattanás utáni viszkózus újramelegítésnek ($T_{\mathrm{reh}} = 1{,}00\times 10^{18}\;\mathrm{GeV}$, termikus leptogenezis $\eta_B^{\mathrm{0\text{-}param}} = 1{,}12\times 10^{-10}$, $\varepsilon_{\mathrm{CP}}^{\mathrm{DI,max}} = 5{,}52\times 10^{-2}$) kovariáns áramlástani elméletével (\texttt{F55 / Gate E31}).
\item \textbf{Ontológiai szintézis és \texttt{Paper 23 (Sub4 v5)} --- Atemporális Morfológiai Test ($M_\tau^\star$), id\H{o}független kvantummechanika, passzív Schur-csillapítás (\texttt{CSCR-T27}) és a 2026-os sötétenergia-válság feloldása:} A pregeometriai $K_4$ Alapmez\H{o} ($B_\tau^{\mathrm{base}}$) és az idempotens Clifford Csíra-projektor ($s_U$, $\Pi_i^2 = \Pi_i$, $r_{\mathrm{seed}} = 1/4$, $4 \to 1 + 3$ felbontás) a $K_B m = j_s$ stacionárius variációs egyensúlyon keresztül egyetlen \emph{atemporális morfológiai válasz-konfigurációt}, a Morfológiai Testet ($[M_\tau^\star] = [\mathcal{F}_\tau(B_\tau^{\mathrm{base}}, s_U, \mathcal{K})]_\sim$) stabilizálja. Mivel $M_\tau^\star$ atemporális, a kvantummechanika id\H{o}független alapképletei ($\hat{H}_*\psi = E\psi$, $[\hat{q},\hat{p}] = \mathrm{i}\hbar$, $\Pi_i^2 = \Pi_i$, $\hat{\mathcal{H}}\Psi = 0$) közvetlenül az $M_\tau^\star$ testet írják le küls\H{o} koordinátaid\H{o} nélkül (\texttt{Sub2 / Paper 22}). A bels\H{o} megfigyel\H{o} az $M_\tau^\star$ testet a kanonikus szül\H{o}i mélység ($\Theta_M$) mentén szeletenként olvassa ki (hullámterjedés és id\H{o}függ\H{o} $\mathrm{i}\hbar\partial_\tau\psi = \hat{H}_*\psi$ Schr\"odinger-egyenlet), míg a rejtett bels\H{o} módusok Feshbach--Schur--Jacobi eliminációja ($C_{\mathrm{FSJ}} = K_0 - H_b = u\,\operatorname{diag}(0,8,0,0) \succeq \mathbf{0}$, \texttt{CSCR-T27}) el\H{o}állítja a passzív terminális Schur-csillapítást, a $\hat{m}_{\mathrm{eff}} = 2{,}12055$ Yukawa-árnyékolást és az $\alpha_{\mathrm{OCSR}} = \mathrm{Tr}[(I_4 - H_b)_{\text{aktív}}] = \frac{5\sqrt{3}}{24}$ szuszceptibilitást. A \textbf{Paper 23 (\texttt{Sub4 v5})} megmutatja, hogy a stacionárius $w_{\mathrm{hor}} = -1$ háttérhorizont szükséges és elégséges a $z_{\mathrm{dip}} \equiv z_{\mathrm{acc}} = 0{,}6316$-os $-6{,}15\%$-os galaktikus völgyhöz ($N=43$ \texttt{JWST/HST} korongon az egységes exponenciális \texttt{COCA}-maggal $\chi^2 = \mathbf{138{,}97}$ vs.\ $160{,}78$, $\Delta\chi^2 = +21{,}80$; algebrai sz\H{u}r\H{o}vel $106{,}42$ vs.\ $121{,}64$), a fal-lassulási átmenet ($q_{\mathrm{wall}} = 0$) illesztése $z_{\mathrm{acc}} = 0{,}6316$-hoz szabad paraméter nélkül $f_{v,0} = \mathbf{0{,}7234}$-ra ($72{,}34\%$) rögzíti az üreg-térfogatarányt, Wiltshire egzakt távolság-integrálja és az $\alpha_{\mathrm{OCSR}} = \frac{5\sqrt{3}}{24}$ súly pedig már kitev\H{o}választás nélkül ($p = 0$, $\Delta k = 0$) látszólagos $(w_0^{\mathrm{app}}, w_a^{\mathrm{app}}) = (\mathbf{-0{,}8885}; \mathbf{-0{,}2782})$ (\texttt{El\H{o}re H2-Sachs}) ill.\ $(\mathbf{-0{,}8843}; \mathbf{-0{,}2042})$ (\texttt{H2a}) értéket ad, és a valódi $N = 1\,580$ \texttt{Pantheon+} szupernóván felülmúlja a standard $\Lambda$CDM-et ($\chi^2_{\mathrm{SN}} = \mathbf{1386{,}21}$ és $\mathbf{1386{,}59}$ vs.\ $\Lambda$CDM $1387{,}33$).
\end{itemize}

\section{Összefoglaló fejl\H{o}dési táblázat}

\begin{center}
\footnotesize
\renewcommand{\arraystretch}{1.25}
\begin{tabular}{>{\raggedright\arraybackslash}p{2.7cm} >{\raggedright\arraybackslash}p{2.1cm} >{\raggedright\arraybackslash}p{3.9cm} >{\raggedright\arraybackslash}p{1.5cm} >{\raggedright\arraybackslash}p{5.0cm}}
\toprule
\textbf{Korszak / Név} & \textbf{Id\H{o}szak} & \textbf{Központi képlet / Struktúra} & \textbf{Paraméter} & \textbf{F\H{o} eredmény és továbblépési ok} \\
\midrule
\texttt{GalaxyTest} / \texttt{SRTM} (\texttt{V1--V5}) & 2026. ápr. 18--23. & Korong-skálahossz ($R_{\mathrm{disk}}$) alapú válasz & $k=2\text{--}3$ & 7 galaxison $\alpha=0{,}405$, de 175 galaxison nagy hiba (\textsc{rmse} $25\text{--}67\;\mathrm{km/s}$). \\
\texttt{MTW} (\texttt{Rev 6}, \texttt{V6}) & 2026. ápr. 23--28. & $F(a_N)=1+\alpha\ln(1+a_0/a_N)$ + \textsc{aqual} Lagrange & $k=1$ ($\alpha=0{,}360$) & 175 \textsc{sparc} galaxison $16{,}62\;\mathrm{km/s}$ ($0{,}97\times$ \textsc{mond}); $\alpha\approx 2a_0/(cH_0)$. \\
\texttt{TPG} + 5D \textsc{rmond}-híd & 2026. ápr. 29. -- máj. 4. & $\int_0^\infty \frac{\alpha}{(1+\tau)^{1+q}}\mathcal{O}(x,\tau)\,d\tau \to$ \texttt{SZ20 \textsc{rmond}} & $k=1$ ($\alpha$) & $c_{\mathrm{mix}}=1/(4+\sqrt{3})$, $c_{\mathrm{GW}}=c$; de Hollinger szuperhalmaz \texttt{18/18 FAIL}. \\
\texttt{HDDA--DTL} & 2026. máj. 5--11. & $\mathcal{A}=(S,\Omega,H)$, $D_x(\mathcal{A})$ + véges kapacitású rétegek & $k=1\text{--}2$ & Környezeti/el\H{o}történeti kapu; 2026.05.08: diszkrét id\H{o}lapok elvetése. \\
\texttt{Tau Core} + \texttt{TPST} + \texttt{Papers 1--19} & 2026. máj. 12. -- aug. & $\Psi(x,\tau) \to 5$ projekciós szektor ($P_{\mathrm{time..gravity}}$) & Szektor-rögzített & \textsc{sparc} zavartsági audit (\texttt{P1}), Higgs $\nu_D=3/10$ (\texttt{P4}), 4D kovariancia (\texttt{P17}). \\
Római Sorozat (\texttt{I--VII-C}) + \texttt{TD} & 2026. aug. -- okt. & $B_\tau^{\mathrm{base}} \to r_{\mathrm{seed}} \to M_\tau^\star$, $\Gamma_{\mathrm{CLP3}}$, 256D $\mathrm{CAR}_7$ & Feltételes csomag & Born/Uhlmann-tétel, Schur-pozitivitás (\texttt{F18}), \texttt{OCSR} koercivitás (\texttt{F16--F20}). \\
\textbf{\texttt{Tau Core v28.1}} ($N=4$ + \texttt{P20--23 / Sub1--Sub4}) & 2026. nov. 12--18. & $N=4 \implies u=\frac{\sqrt{3}}{24}$ + $M_\tau^\star$ + $\mathcal{S}_{\mathrm{Tau}}$ + PG áramlástan + üreg-órák ($f_{v,0}=72{,}34\%$) & \textbf{$k = 0$} & $\alpha_{\mathrm{OCSR}}=5u$, $z_{\mathrm{dip}}=0{,}6316$, \texttt{Pantheon+} $p=0$ $\chi^2=1386{,}21 < 1387{,}33$, \textbf{35/35 PASS}. \\
\bottomrule
\end{tabular}
\end{center}

\end{document}
